\documentclass[10pt,a4paper]{amsart}
\usepackage[margin=19mm]{geometry}
\usepackage{amsmath,amssymb,mathtools}
\usepackage[scr=rsfs,cal=euler]{mathalfa}
\usepackage{xfrac}
\usepackage[unicode,hidelinks,bookmarks=false]{hyperref}
\newcommand{\ie}{i.e.\ }
\newcommand{\cf}{cf.\ }
\newcommand{\resp}{respectively\ }
\newcommand{\Subsection}{\subsection}
\providecommand{\nobreakdash}{\nobreak\hbox{-}}
\setlength{\emergencystretch}{2em}
\multlinegap0pt
\usepackage{bm}
\usepackage{bbm}
\usepackage{dsfont}
\usepackage{xspace}
\usepackage{cancel}
\usepackage{mathtools}
\usepackage{amsthm}
\makeatletter
\@ifundefined{theorem}{\newtheorem{theorem}{Theorem}}{}
\@ifundefined{proposition}{\newtheorem{proposition}[theorem]{Proposition}}{}
\makeatother
\title[Large induced subgraphs with prescribed degree parity]{Large induced subgraphs with prescribed degree parity}
\author{Jiangdong Ai, Qiwen Guo, Gregory Gutin, Yiming Hao, Anders Yeo}
\date{Source: arXiv:2509.01428v2 (CC BY 4.0)}
\begin{document}
\maketitle

\noindent\emph{Source and licence.} Large induced subgraphs with prescribed degree parity. Version arXiv:2509.01428v2 (CC BY 4.0), distributed under the Creative Commons Attribution 4.0 International licence, as recorded in the arXiv licence metadata for this exact version and shown on the abstract page https://arxiv.org/abs/2509.01428v2. Full rights notice: licenses/arxiv-2509.01428-CC-BY-4.0.txt.\par\smallskip

\noindent\textit{Editorial scope.} Counted unit: the Proposition whose source label is \texttt{None} in the article (source0.tex lines 606--611, source label \texttt{None}), delivered together with the article's own complete original proof of it (source0.tex lines 612--724, one \texttt{proof} environment of 113 lines). No mathematics is written, repaired or re-derived: statement and proof are the article's own text line for line. Required context retained uncounted: none. Every definition, display and label of those lines is retained unchanged. Omitted from this package and not redistributed: 1--605, 725--839. Nothing inside a retained excerpt is omitted or altered: every retained body is delivered exactly as the article's lines are written, so no omission receipt of any kind accompanies this package. Citations: the retained excerpts carry no citation command at all, so no bibliography entry of the article is transcribed and no reference is redistributed. Reference adaptation: none was needed and none was made; every reference inside the delivered unit resolves inside it, so no unresolved reference or citation is produced. Printed slips kept literal: no printed word, symbol, hypothesis or number of the counted statement or of its proof was altered. Layout and class: the article's own class is \texttt{article} with packages including amsmath, amssymb, amscd, amsthm, amsfonts, graphicx, mathrsfs, hyperref, color, float, epstopdf, tikz, caption, subcaption; the wrapper is the lane's pinned \texttt{amsart} editorial wrapper at 10pt on A4, which restates the theorem-like environments the retained text needs and adds the article's own macro definitions that the retained text needs (none), with extra wrapper packages (bm, bbm, dsfont, xspace, cancel, mathtools). The delivered numbering is the unit's own. Assets: no retained span contains a figure environment, a graphics-inclusion command, a PDF-page inclusion, a table or a TikZ picture; no image file is copied into this package, the package declares no asset dependency and no image file is redistributed with it. Licence: version arXiv:2509.01428v2 of this article is distributed under the Creative Commons Attribution 4.0 International licence, as recorded in the arXiv licence metadata for this exact version and shown on the abstract page https://arxiv.org/abs/2509.01428v2. A full rights notice is delivered at licenses/arxiv-2509.01428-CC-BY-4.0.txt. Characters: every retained excerpt is pure ASCII, so the delivered engine, XeLaTeX with its default Latin Modern text font, reports no missing character.
\begin{proposition}
For the star \(K_{1,m}\), where \(m\ge 1\),
\[
f_{oe}(K_{1,m})=\left\lfloor \frac{m+1}{2}\right\rfloor .
\]
\end{proposition}

\begin{proof}
Let \(c\) be the center of \(K_{1,m}\), and let \(L\) be the set of its \(m\) leaves.
Fix a labeling
$
\ell:V(K_{1,m})\to\{0,1\}
$ and 
let
$
a=|L\cap \ell^{-1}(0)|, b=|L\cap \ell^{-1}(1)|,
$
so that \(a+b=m\).

We first determine \(h_\ell(K_{1,m})\). Let \(U\subseteq V(K_{1,m})\) be
\(\ell\)-admissible.

If \(c\notin U\), then every leaf in $U$ has degree \(0\) in \(K_{1,m}[U]\).
Thus, $ |U|\le a.$
Conversely, the set of all \(0\)-labeled leaves is \(\ell\)-admissible, so this
case the maximum of  \(|U|\) is $a$.

Now suppose \(c\in U\). Since every leaf in $U$ has degree \(1\) in
\(K_{1,m}[U]\), every such leaf must be \(1\)-labeled. If \(t\) denotes the
number of such leaves, then the center has degree \(t\) in \(K_{1,m}[U]\),
so we also require
$
t\equiv \ell(c)\pmod 2.
$
Thus the maximum of $|U|$ is
\(
1+t_\ell,
\)
where \(t_\ell\) is the largest integer satisfying
\(
0\le t_\ell\le b, t_\ell\equiv \ell(c)\pmod 2,
\)
if such an integer exists. If no such integer exists, then no \(\ell\)-admissible
induced subgraph containing the center exists. The only exceptional case is
\(b=0\) and \(\ell(c)=1\). Hence
\[
h_\ell(K_{1,m})=\max\{a,\ 1+t_\ell\},
\]
where the second term equals \(0\) when \(b=0\) and \(\ell(c)=1\).

We now prove the lower bound $h_\ell(K_{1,m})\ge \left\lfloor\frac{m+1}{2}\right\rfloor.$ If
\(
a\ge \left\lfloor\frac{m+1}{2}\right\rfloor,
\)
then the set of all \(0\)-labeled leaves is already an \(\ell\)-admissible set
of size at least \(\left\lfloor\frac{m+1}{2}\right\rfloor\).

Otherwise,
\(
a<\left\lfloor\frac{m+1}{2}\right\rfloor.
\)
Since \(a+b=m\), this implies
\(
b\ge \left\lceil\frac{m+1}{2}\right\rceil .
\)
Among the two largest integers \(b\) and \(b-1\), one has parity \(\ell(c)\).
Therefore there exists an integer \(t\le b\) with
\(
t\equiv \ell(c)\pmod 2
\)
and
\(
t\ge b-1.
\)
Thus the maximum size of $U$ containing \(c\) is at least
\[
1+(b-1)=b\ge \left\lceil\frac{m+1}{2}\right\rceil
\ge \left\lfloor\frac{m+1}{2}\right\rfloor .
\]
Therefore, for every labeling \(\ell\),
\[
h_\ell(K_{1,m})\ge \left\lfloor\frac{m+1}{2}\right\rfloor .
\]
Hence
\(
f_{oe}(K_{1,m})\ge \left\lfloor\frac{m+1}{2}\right\rfloor .
\)

It remains to prove the matching upper bound. We exhibit labelings for which no larger
\(\ell\)-admissible induced subgraph exists.
First suppose \(m=2r\). Label \(r\) leaves by \(0\) and \(r\) leaves by
\(1\). Choose the label of the center so that
\(
\ell(c)\not\equiv r\pmod 2.
\)
Then an \(\ell\)-admissible set not containing the center has size at most \(r\), while
an \(\ell\)-admissible set containing the center can contain at most \(r-1\) of the
\(1\)-labeled leaves, and therefore has size at most \(1+(r-1)=r\). Hence
\[
h_\ell(K_{1,2r})\le r=\left\lfloor\frac{2r+1}{2}\right\rfloor .
\]

Now suppose \(m=2r+1\). Label \(r+1\) leaves by \(0\) and \(r\) leaves
by \(1\). Choose the label of the center so that
$
\ell(c)\not\equiv r\pmod 2.
$
Then an \(\ell\)-admissible set not containing the center has size at most \(r+1\),
while an \(\ell\)-admissible set containing the center can contain at most \(r-1\)
\(1\)-labeled leaves, and therefore has size at most \(r\). Hence
\[
h_\ell(K_{1,2r+1})\le r+1
=
\left\lfloor\frac{2r+2}{2}\right\rfloor .
\]
Combining the lower and upper bounds gives
\[
f_{oe}(K_{1,m})=\left\lfloor \frac{m+1}{2}\right\rfloor .
\]
\end{proof}

\end{document}
